\section{Problema dell'ordinamento}
Il problema dell'ordinamento dei dati in una struttura è uno dei più comuni che si incontra nella pratica e consente di introdurre vari strumenti di analisi e tecniche di progettazione standard. Formalmente viene definito come: \begin{itemize}
	\item \textbf{Input}: Sequenza di $n$ oggetti $\left<a_1, ..., a_n\right>$ sulla quale è definita una relazione di ordine totale.
	\item \textbf{Output}: Permutazione $\left<a_1', ..., a_n'\right>$ della sequenza di input tale che $a_1' \leq ... \leq a_n'$.
\end{itemize}
\noindent In pratica, data per esempio una sequenza di numeri in input $\left<1, 4, 2, 7, 5, 3\right>$, l'algoritmo risolutivo ritornerà in output $\left<1, 2, 3, 4, 5, 7\right>$. La sequenza specifica in entrata è detta \textbf{istanza} del problema e per essere elaborata è necessario che appartenga al dominio per cui l'algoritmo è definito.\par
Diciamo che un algoritmo è \textbf{corretto} se, per ogni istanza valida del problema, termina e produce un output conforme alla specifica e sebbene si tratti di un concetto intuitivamente semplice, esistono molti algoritmi adatti a questo compito, la cui scelta è determinata da alcuni fattori come la \textbf{dimensione} del numero di dati, il loro \textbf{ordinamento iniziale}, eventuali \textbf{vincoli} o anche limiti fisici dell'elaboratore. 

%

\section{Insertion Sort}
Algoritmo efficiente per ordinare un piccolo numero di elementi, opera allo stesso modo in cui si ordinerebbe un mazzo di carte. Dato in input un array $A$, mantengo una parte ordinata ed una non ordinata; ad ogni iterazione si prende un elemento dalla seconda e lo si inserisce nella posizione corretta, trovata confrontando tutte le carte a partire dall'inizio. \begin{itemize}
	\item \textbf{Caso pessimo}: $\Theta(n^2)$
	\item \textbf{Caso ottimo}: $\Theta(n)$. Si ottiene quando l'array è già ordinato.
	\item \textbf{Caso medio}: $\Theta(n^2)$
	\item \textbf{Ordinamento in loco}: Sì
	\item \textbf{Stabilità}: Sì
\end{itemize}
\begin{algorithm}[H]
	\caption{InsertionSort(A,n)}
	\begin{algorithmic}[1]
		\Require $A$: Array, $n$: Dimensione array
		\For{$i \gets 2 \textbf{ to } n$} \Comment{Scorre da startPos+1=2 fino a n dimensione}
		\State $\text{key} \gets A[i]$
		\State $j \gets i-1$ \Comment{L'indice j è usato per la sezione ordinata. Valori: [1,n]}
		\State
		\While{$j > 0 \textbf{ AND } A[j] > \text{ key}$} \Comment{Non sei all'inizio e key minore di A[i]}
		\State $A[j+1] \gets A[j]$ \Comment{Sposta a destra di una posizione il valore maggiore di key}
		\State $j \gets j-1$ \Comment{Spostati a sinistra di una posizione}
		\EndWhile
		\State $A[j+1] \gets key$ \Comment{Inserisci key nella posizione corretta}
		\EndFor
	\end{algorithmic}
\end{algorithm}
\noindent La valutazione della complessità di questo algoritmo vede analizzati due elementi principali: il ciclo \textbf{for}, avente tre istruzioni (gli assegnamenti) di complessità costante, e il ciclo \textbf{while} interno al primo, il quale, potendo scorrere potenzialmente l'intero array ad ogni iterazione, aumenta notevolmente i tempi di esecuzione.\par
Dove richiede un tempo di lavoro generalmente alto rispetto ad altri algoritmi, la sua potenzialità sta nel risparmio della memoria; infatti richiede solo tre variabili per essere eseguito e non richiede alcun supporto esterno per la gestione dell'array. Diciamo, per questa ragione, che l'algoritmo \textbf{opera in loco}. Un'altra sua caratteristica è che gli elementi uguali non vengono riordinati fra di loro poiché il confronto effettuato è stretto. In queste condizioni, diciamo che l'algoritmo è \textbf{stabile}.

%

\section{Merge Sort}
Algoritmo basato sulla tecnica del \textbf{Divide et Impera}, ovvero nella suddivisione del problema in sottoproblemi della stessa natura di dimensione via via più piccola.\par
Diviso l'array di partenza in due parti uguali, si ordinano entrambe applicando ricorsivamente l'algoritmo, per poi combinare queste due nuove sezioni con un sottoalgoritmo \textbf{merge} di supporto, prendendo volta per volta il valore più piccolo contenuto in entrambi, fino ad esaurire tutti i valori. \begin{itemize}
	\item \textbf{Caso pessimo}: $\Theta(nlog(n))$
	\item \textbf{Caso ottimo}: $\Theta(nlog(n))$
	\item \textbf{Caso medio}: $\Theta(nlog(n))$
	\item \textbf{Ordinamento in loco}: No, richiede memoria ausiliaria.
	\item \textbf{Stabilità}: Sì
\end{itemize}
\begin{algorithm}[H]
	\caption{merge(A, left, center, right)}
	\begin{algorithmic}[1]
		\Require A: Array, left: Posizione iniziale, center: Centro, right: Posizione finale
		\State $i\gets \text{left}$
		\State $j\gets \text{center}+1$
		\State $k = 1$ \Comment{Indice dell'array di supporto}
		\State $B[right-left+1]$ \Comment{Array di supporto di dimensione right-left+1}
		\State
		\While{$i \leq \text{center}$ \textbf{AND} $j\leq\text{right}$} \Comment{Fintanto che i sottarray non sono finiti}
		\If{$A[i] \leq A[j]$} \Comment{Se l'elemento del primo è minore, copia quello e avanza}
		\State $B[k]\gets A[i]$
		\State $i\gets i+1$
		\State $k\gets k+1$
		\Else \Comment{Altrimenti procedimento analogo, ma per il secondo array}
		\State $B[k]\gets A[j]$
		\State $j\gets j+1$
		\State $k\gets k+1$ 
		\EndIf
		\EndWhile
		\While{$i \leq \text{center}$} \Comment{Copia il primo array ordinato}
		\State $B[k]\gets A[i]$
		\State $i\gets i+1$
		\State $k\gets k+1$
		\EndWhile
		\While{$j\leq \text{right}$} \Comment{Copia il secondo array ordinato}
		\State $B[k]\gets A[j]$
		\State $j\gets j+1$
		\State $k\gets k+1$
		\EndWhile
		\For{$k\gets \text{left to right}$} \Comment{Sovrascrivi l'array di base con quello ordinato}
		\State $a[k]\gets b[k-\text{left}]$
		\EndFor
	\end{algorithmic}
\end{algorithm}
\begin{algorithm}[H]
	\caption{mergeSort(A, left, right)}
	\begin{algorithmic}[1]
		\Require A: Array, left: Posizione iniziale, right: Posizione finale
		\If{left $<$ right}
		\State center $=$ (left$+$right)$/2$
		\State mergeSort(A, left, center) \Comment{Ordina la parte prima della metà}
		\State mergeSort(A, center$+1$, right) \Comment{Ordina la parte dopo la metà}
		\State merge(A, left, center, right) \Comment{Fondi i due sottoarray ordinati}
		\EndIf
	\end{algorithmic}
\end{algorithm}
\noindent Diversamente rispetto ad Insertion Sort, per valutare la complessità di questo algoritmo bisogna considerare la ricorsione. Sapendo che ad ogni passo ricorsivo il problema è diviso in \textbf{due sottoproblemi} $(n/2)$, dove si effettuano \textbf{due chiamate ricorsive} $(2T)$, e che il merge scorre tutti gli elementi \textbf{almeno una volta} $\Theta(n)$, possiamo dire che Merge Sort è rappresentato da: \[T(n) = 2T(n/2) + \Theta(n)\]
\noindent Possiamo usare il teorema dell'esperto per la dimostrazione. Convertendo la ricorrenza alla forma generale, otteniamo: \[a=2,\quad b=2,\quad f(n) = \Theta(n),\quad log_2(2)=1\implies n^1=n\]
\noindent Confrontando $f(n)$ con la spartiacque, notiamo che appartengono allo stesso ordine di grandezza, facendoci trovare nel caso 2 del teorema, infatti: \[f(n) = \Theta\left(n^{log_b(a)}\right)\implies T(n) = \Theta(nlog(n))\]

%

\section{Struttura Heap e Heap Sort}
Introduciamo ora la prima struttura di dati del corso. Nelle sezioni precedenti è stato menzionato il concetto di \textbf{albero binario}, una struttura composta in cima da un nodo, detto \textbf{radice}, il quale, in quanto binario, si dirama in due direzioni verso il basso attraverso \textbf{archi}, creando due nodi figli. In particolare, quando un nodo figlio è alla base dell'albero si dice \textbf{foglia}.\par
Come è possibile intuire, gli alberi binari sono organizzati in livelli, i quali indicano la \textbf{profondità}, ovvero la distanza dalla radice. Quando un livello ha il numero massimo di nodi ivi ottenibili si dice \textbf{completo}. Tuttavia possono esistere alberi \textbf{semi-completi}, quando la base non è completa.\par
Un vantaggio di questa struttura è la semplicità della ricerca dei nodi; infatti, la dinamica si dispiega in tre casi: \begin{itemize}
	\item \textbf{Per figlio sinistro}: $\text{indice}\cdot 2+0$
	\item \textbf{Per figlio destro}: $\text{indice}\cdot 2+1$
	\item \textbf{Da figlio a padre}: $\text{indice}/2$
\end{itemize}
\begin{minipage}[H]{0.4\linewidth}
	\begin{center}
		\begin{tikzpicture}[
			level distance=1.5cm,
			level 1/.style={sibling distance=4cm},
			level 2/.style={sibling distance=1.3cm}
			]
			
			\node {$1$}
			child {node {$2$}
				child {node {$4$}}
				child {node {$5$}}
			}
			child {node {$3$}
				child {node {$6$}}
				child {node {$7$}}
			};
		\end{tikzpicture}
		Albero binario completo
	\end{center}
\end{minipage}
\hspace{0.2\linewidth}
\begin{minipage}[H]{0.4\linewidth}
	\begin{center}
		\begin{tikzpicture}[
			level distance=1.5cm,
			level 1/.style={sibling distance=4cm},
			level 2/.style={sibling distance=1.3cm}
			]
			
			\node {$1$}
			child {node {$2$}
				child {node {$4$}}
				child {node {$5$}}
			}
			child {node {$3$}
				child {node {$6$}}
			};
		\end{tikzpicture}
		Albero binario semi-completo
	\end{center}
\end{minipage}
\vspace{0.01\linewidth}

\noindent Un \textbf{heap} è una struttura dati implementata tramite un array, considerata sotto forma di un albero binario semi-completo, i cui nodi rappresentano gli elementi al suo interno. Deve inoltre rispettare la proprietà che lo contraddistingue: in un \textbf{MAX-Heap} ogni nodo è maggiore o uguale ai propri figli. Analogamente, è possibile definire un \textbf{MIN-Heap} dove vale il contrario.
\noindent Nella gestione di questa struttura dati, tuttavia, in necessità di inserimento nodi, non necessariamente si potrà rispettare sempre la proprietà, ragion per cui utilizziamo un algoritmo \textbf{heapify()} ad-hoc, con lo scopo di prendere questo problema locale e trasferirlo alle foglie, dove sarà risolto.
\begin{algorithm}[H]
	\caption{heapify(A, i)}
	\begin{algorithmic}[1]
		\Require $A$: Array, $i$: Indice
		\State $l\gets left(i)$, $r\gets right(i)$ \Comment{l prende figlio-sx, r prende figlio-dx}
		\State
		\If{$l\leq A.heapSize$ \textbf{AND} $A[l] > A[i]$} \Comment{Se figlio-sx ha valore $>$ i, non è heap.}
		\State largest $\gets l$
		\Else 
		\State largest $\gets i$ \Comment{Altrimenti segna che il più grande è i, come inteso.}
		\EndIf
		\State
		\If{$r \leq \text{A.heapSize}$ \textbf{AND} $A[r] > A[\text{largest}]$} \Comment{Controllo figlio destro}
		\State largest $\gets r$
		\EndIf
		\State
		\If{largest $\neq i$} \Comment{Se la proprietà non è valida}
		\State switch(A[i], A[largest]) \Comment{Scambia la posizione dei due nodi, largest è l'anomalia}
		\State heapify(A, largest) \Comment{Sposta ricorsivamente largest finché non rispetta la regola}
		\EndIf
	\end{algorithmic}
\end{algorithm}
\noindent La complessità di questo algoritmo è equivalente alla sua profondità, quindi il caso peggiore scende lungo l'altezza dell'albero fino alle foglie, dando come tempo di esecuzione $T(n) \in O(log(n))$.



% TODO Riprendi da buildHeap, ma prima aggiungi la dimostrazione della ricorrenza a pag.149



\begin{comment}	
	-- è possibile passare da array a heap?
	
	buildHeap(A) {
		A.heapSize = A.length;
		
		// Tutti i nodi con chiave maggiore di i sono a posto. Chiama heapify, se la proprietà della heap è ok, va bene, altrimenti la riordina come vista nell'algoritmi prima.
		for (i = A.length/2; i > 0; i--) {
			heapify(A, i);
		}
	}
	
	28-10-25 Algoritmi
	
	Formula per il totale dei nodi (somma totale elementi in una serie geometrica): \frac{1-2^{h+1}}{1-2} = 2^{h+1}-1
	Al livello più basso abbiamo in totale n/2 nodi, con 0 operazioni. 		// n/2 perché consideri solo quel livello senza guardare i superiori, quindi non è n.
	Al precedente, n/4, con 1 operazione.
	Al precedente ancora, n/8, con 2 operazioni. e così via.
	
	Voglio una formula che dà la complessità dell'algoritmo
	\sum_{i=0} i\frac{n}{2^{i+1}}		// i sono le operazioni fatte in quel livello, nel denominatore comporrà il totale dei nodi.
	
	\sum_{i=0} i\frac{n}{2^{i+1}} \leq n\sum_{i=0}^{+\infty} i(\frac{1}{2})^{i+1}.		// La sommatoria converge ad un numero finito, quindi è lineare.
	
	// Vediamo come converge
	\sum_{i=0}^{\infty} iq^{i+1}		// Somiglia alla sommatoria delle derivate di una serie geometrica.
	&= \sum_{i=1}^{\infty} iq^{i+1}
	&= q^2 \sum_{i=1}^{\infty} iq^{i-1}		// Sommatoria delle derivate dei termini di una serie geometrica. Questa converge alla derivata della somma. E' un numero finito.
	
	Esce che la sommatoria è una costante e la funzione è lineare. Figo, hai le basi per fare un algoritmo di ordinamento basato sull'heap.
	
	--- HeapSort
	HeapSort(A) {
		buildHeap(A);		// Crea l'heap
		
		for (i = A.length; i > 1; i--) {
			switch(A[1], A[i]);
			A.heapSize--;					// Riduco la dimensione dell'heap
			heapify(A,1);
		}
	}
	
	Perché funziona?
	Prendi l'array già organizzato come heap. L'elemento più piccolo è la radice, quello più grande, in fondo.
	Scambio l'elemento in fondo con il primo e poi decremento la dimensione. Prendo l'elemento più grande e lo metto in cima, poi con heapify ordina l'array fin quando non scorre tutti
	gli elementi.
	
	- Complessità: buildHeap è lineare, switch è costante, heapify è logn. il corpo del ciclo for nel peggiore dei casi è nlogn.
	limite superiore nlogn, il limite inferiore è log(n!) = \theta(logn^n) = \theta(nlogn)
	- Ordina in loco, perché ogni operazione è fatta sullo stesso array
	- è stabile? No. Prova a usare un array con tutti gli elementi uguali.
	
	RICORDARE n! = n^n.
\end{comment}

%

\section{Quick Sort classico e probabilistico}





\begin{comment}
	--- QuickSort
	Ha di base divide et impera come nel mergeSort.
	Prendiamo l'array, lo divisiamo in due parti tali che tutti gli elementi della prima parte siano minori o uguali di quelli della seconda.
	Poi ordino la parte di sinistra e dopo la parte di destra.
	
	QuickSort(A, P, r) {		// P è dove parte l'array, r è dove finisce l'array.
		if (P < r) {
			q = partition(A,P,r);	// Divido l'array secondo la logica su scritta. q è l'elemento perno, pivot.
			QuickSort(A,P,q);		// Ordino la sinistra
			QuickSort(A,q+1,r);		// Ordino la destra
		}
	}
	
	// Complessità: proporzionale agli spostamenti di i e j. è lineare.
	partition(A, P, r) {
		x = A[P];			// Prende il primo elemento e lo rende quello medio per spostare gli elementi.
		i = P-1;
		j = r+1;
		
		while (1) {
			while (A[j] > x) {
				j = j-1;
			}
			
			while (A[i] < x) {
				i = i+1;
			}
			
			if (i < j) {
				switch(a[i], A[j]);
			} else {
				return j;
			}
		}
	}
	
	03-11-25 Algoritmi
	
	--- QuickSort (A,P,r)
	
	QuickSort (A,P,r) {
		if (P<r) {
			q = partition(A,P,r);
			
			QuickSort(A,P,q);
			QuickSort(A,q+1,r);
		}
	}
	
	partition (A,P,r) {	// Complessità: n
		x = A[P];
		i = P-1;
		k = r+1;
		
		while (1) {
			while (A[j]<=x) {		// Sposta finché non vale più la regola
				j--;
			}
			
			while (A[i]>=x) {		// Sposta finché non vale più la regola
				i++;
			}
			
			if (i < j) {			// Sei qua solo se hai violato ambo i casi
				switch(A[i], A[j]);	// Cambia di posizione
			} else {
				return j;			// Torna l'indice j.
			}
		}
	}
	
	T(n) = \theta(n) + T(n-1) + T(1) \equiv \theta(n) + T(n-1)
	&= n + (n-1) + T(n-2)
	&= n + (n-1) + (n-2) + T(n-3)
	&= \sum_{i=1}^n i = \frac{n(n+1)}{2}		// Il caso pessimo di QuickSort è quadratico, O(n^2).
	
	Per risolvere pseudo-completamente il problema del caso pessimo (array ordinato) si può scegliere un elemento perno a cazzo oppure mescolare
	l'array in toto.
	
	randPartition (A,P,r) {
		i = random(P,r);		// Scegli un numero a caso compreso fra P ed r.
		switch(A[P], A[i]);
		
		return partition(A,P,r);
	}
	
	Qua parliamo di caso medio perché c'è la probabilità in gioco. Infatti hai una prob 1/n di pescare un dato elemento perno.
	Serva una funzione che associa ad ogni n una variabile aleatoria.
	
	Se vedo T(n) come la funzione di variabile aleatoria, abbiamo:
	T(n) &= \theta(n) + 1/n(T(1) + T(n-1)) + 1/n(T(1) + T(n-1)) + 1/n(T(2) + T(n-2)) + ... + 1/n(T(n-1) + T(1))\\
	&= \theta(n) + 1/n(T(1) + T(n-1)) + \sum_{i=1}^{n-1} 1/n(T(i) + T(n-1))
	&= \theta(n) + \sum_{i=1}^{n-1} 1/n(T(i) + T(n-1)) // Questo perché T(1)è costante, fanculo. Poi il caso pessimo n^2, se moltiplicato a 1/n
	// ci dona il valore n, che sta sotto \theta(n). Quindi si può rimuovere la componente
	&= \theta(n) + 2/n \sum_{i=1}^{n-1} T(i)
	
	// TODO vedere dimostrazione con metodo di sostituzione, dimostra che il caso medio è nlogn
	
	// Question: Hai 12 sfere, tutte uguali. (ci puoi fare quel che vuoi). Fra d di loro ce ne potrebbe essere una con un peso diverso dalle altre.
	// Abbiamo una bilancia coi due piatti. Hai tre pesate per dire se c'è una sfera diversa dalle altre e se pesa di più o meno.
	
	Esiste un criterio per dire se esiste un algoritmo per risolvere un determinato problema.
	Con quello delle sfere appena visto, abbiamo tre possibili risultati; palla leggera, palla pesante, palle uguali.
	Con tre pesate massimo per un tot di palle.
	
	La prima è obbligata. 4 palle vs 4 palle. 3pesate*3scenari = 9.
	
	// QuickSort iterativo
	QuickSort (A,P,r) {
		while (P<r) {
			q = partition(A,P,r);
			
			QuickSort(A,P,q);
			p = q+1;
		}
	}
\end{comment}

%

\section{Counting Sort}





\begin{comment}
	04-11-25 Algoritmi
	Soluzione palle:
	Prima pesata: 4 palle per piatto. (1,2,3,4) - (5,6,7,8)
	Se il piatto tilta, o c'è una palla che pesa di più o di meno, altrimenti dovrai prendere il terzo gruppo di palle restanti.
	
	Seconda pesata: (1,5,7) - (2,6,9)
	Se il piatto tilta con sx più pesante possiamo avere che 5+ o 7+ o 2-.
	Se il piatto tilta con dx più pesante possiamo avere che 1- o 6+.
	Se i piatti sono uguali, possiamo avere 3- o 4- o 8+.
	
	Terza pesata:
	
	Possiamo vedere un algoritmo che lavora per confronti come un albero di decisione, che ha radice come il totale delle sue permutazioni [n!] e termina quando si 
	ottiene la singola istanza, la foglia. La profondità è ugualmente data dal logaritmo. Nel nostro caso specifico: log(n!) = log(n^n) = nlogn.
	Un algoritmo che lavora al meglio partiziona l'insieme delle soluzioni e non ci devono essere soluzioni ripetute. Non può fare meglio di così, ammesso che esista.
	Questo è il limite inferiore.
	Se esiste, possiamo provare il limite superiore, se non esiste abbiamo almeno questo limite inferiore teorico.
	
	--- Algoritmi di ordinamento lineari
	Se avessimo un'informazione in più per l'ordinamento che consente di riordinare saltando dei passaggi? A questo punto potremmo scendere ulteriormente in execTime.
	Quindi sotto nlogn.
	
	Supponiamo di avere un array di int, tutti compresi fra 1 e k, dove k è arbitrario. Posso usare questa k per ordinare più velocemente.
	Contando gli elementi è in grado di capire la posizione degli elementi.
	
	// Vuole ordinare gli elementi di A, compresi fra 1 e K, per poi metterli in B.
	CountingSort(A, B, K)				// Array A, Array B, int K.
	for (i<-1 to K) C[i] <- 0		// Questa riga ha complessità K
	for (j<-1 to length[A])			// length[A] è la lunghezza di A.	Questo ciclo ha compl. n
	C[A[j]]++					// L'indice è visto come l'elemento A[j], il quale è incrementato.
	for (i<- 2 to K)				// Questo for ha compl. K
	C[i] <- C[i] + C[i-1]		// Elementi con chiave minore di i
	for (j <- length[A] down to 1)	// Scorri l'array alla rovescia fino al primo elemento (non è 0 qui, in pseudo è 1.)	Il for qui ha compl. n.
	B[C[A[j]]] <- A[j]
	C[A[j]]--					// Questa linea garantisce stabilità. Decrementa l'indice dei numeri contati, sicché si possano inserire alla posizione
	// antecedente in B.
	
	La complessità dell'algoritmo è n+k.
	
	A = {3, 6, 4, 1, 3, 4, 1, 4}
	B = {1, 1, 3, 3, 4, 4, 4, 6}		// Questa è la risposta
	// C al secondo for
	C = {0, 0, 0, 0, 0, 0} \implies {0, 0, 1, 0, 0, 0} \implies {0, 0, 1, 0, 0, 1} \implies 0 0 1 1 0 1 \implies 1 0 1 1 0 1 \implies 1 0 2 1 0 1 \implies ...
	// C dice quindi per ogni indice dice quanti elementi ci sono con quel valore. Finalmente, \implies {2, 0, 2, 3, 0, 1}
	
	// C al terzo for
	C = {2, 2, 4, 7, 7, 8}
	
	Potrebbe ordinare anche elementi da 1 a n, basta mettere k=n.
	Essendo che puoi porre k=n, il tempo di esecuzione dipende da n. Se è n^2 è quadratico, per esempio.
	Qui le costanti hanno un impatto rilevante. Con un k molto grande, tipo 2^{63}, è impraticabile, ma teoricamente cresce sempre linearmente.
\end{comment}

%

\section{Radix Sort}





\begin{comment}
	--- RadixSort
	
	329, 457, 657, 839, 436, 720, 355.
	
	Ordina a partire dalla cifra meno significativa tramite counting sort. Otteniamo al primo ciclo:
	720, 355, 436, 457, 657, 329, 839.
	
	Al secondo, numeri ordinati per le due cifre meno significative:
	720, 329, 436, 839, 355, 457, 657.
	
	Al terzo, numeri ordinati per le tre cifre meno significative, quindi tutte.:
	329, 355, 436, 457, 657, 720, 839.
	
	n numero di elementi
	k la base
	l il numero di cifre
	
	Complessità: l(n+k).
\end{comment}

%

\section{Bucket Sort}





\begin{comment}
	--- BucketSort
	Prendi lo spazio degli elelìmenti e dividilo in n bucket equiprobabili.
	Assegna i vari oggetti al bucket rilevante. Poi ordina i bucket e concatena i risultati.
	
	Ci si aspetta che in media ci sia un elemento per bucket.
	In tempo lineare ordina i bucket, poi concatenando hai finito.
	
	
	
	10-11-25 Algoritmi
	
	// Ordinamento lessicografico - ordina caratteri di stringhe
	BucketSort, assegna ogni elemento dell'array al suo bucket, ordina i singoli buckets e hai l'array ordinato.
	
	Definizione formale
	Suppongo VA X_{ij} = \begin{cases}
		1 se A[i]\in Bucket_j\\
		0 altrimenti
	\end{cases}. Capiamo che la VA è di Bernoulli. Un valore ha probabilità 1/n di essere in un bucket, in quanto sono equiprobabili. Le variabili sono anche indipendenti fra di loro.
	Quindi P[x_{ij}= 1] = 1/n = p.
	
	X_{j} = \sum_i X_{ij}.		// Conta il numero di elementi che finiscono nel bucket j, quindi quelli che valgono 1.
	
	La complssità è data quindi da \sum_j X_j^2		// La complessità è quindi data dalla somma dei bucket j che hanno valore 1, gli altri non sono toccati.
	La complessità media è invece data dal valore atteso: \sum_j E[X_j^2]
	
	E[X_{ij}] = 1/n
	Var[X_{ij}] = p*q \implies 1/n(1 - 1/n)
	E[X_j] = \sum_i E[X_{ij}] = n(1/n) = 1		// Mediamente in ogni bucket m'aspetto un elemento.
	Var[X_j] = \sum_i Var[X_{ij}] = n(1/n)(1 - 1/n) = 1 - 1/n
	E[X_j^2] = Var[X_j] + E^2[X_j] = 1 - 1/n + 1^2 = 2 - 1/n
	
	Dunque: \sum_j E[X_j^2] \implies n(2 - 1/n) = 2n - 1
	
	r = \sqrt(x^2 + y^2)
	\theta = arctan(y/x)
	
	Elementi distribuiti uniformemente in un cerchio di raggio r. Ordinare gli elementi in fase da 0 a 2\pi.
	Bucket = 2\pi / n oppure raguionando per angoli abbiamo \alpha_i = i * 2\pi/n
	
	E se volessi ordinare gli elementi per modulo crescente?
	|r| = r*\sqrt(i/n)
\end{comment}

%

\section{Problema della selezione}





\begin{comment}
	--- Problema della selezione
	Dato un array di oggetti sulla quale è definita una relazione di ordinamento. Trovare il valore in posizione i.
	
	Interfaccia data da Select(A,i)
	Meno di n-1 confronti non posso farne; ma in tal caso, se voglio un max o un min, avremo verosimilmente due valori vincitori.
	\Omega(n)
	
	// Tip: divisione in tre parti dell'array, a mò di quickSort
	
	
	
	
	
	11-11-25 - Algoritmi
	
	Per trovare min e max, la complesità è 2n-2, perché esegui due volte uno stesso algoritmo di complessità n-1.
	
	-- Ricerca contemporanea di minimo e massimo
	Dato un array, ogni elemento potrebbe essere potenziale min e potenziale Max, li segnamo rispettivamente come m e M.
	Confrontando elementi possiamo escludere determinati elementi dalla loro potenzialità.
	
	Prendiamo la somma m+M, la quale deve essere uguale a 2 quando la ricerca è conclusa. Questo perché m ed M rappresentano ogni elemento dell'array che è potenzialmente min o max.
	Con i confronti è probabile che:
	- Il valore non diminuisca (Confronto inutile)
	- Il valore diminuisca di 1 (Un valore perde la potenzialità in ambito di due confronti)
	- Il valore diminuisca di 2 (Due valori perdono la potenzialità in ambito di due confronti)
	
	Generalmente, al'inizio: m+M = 2n.
	Dopo x volte: m+M = 2n-2x
	
	Per n pari abbiamo che:
	Per n/2 volte, calo di 2.
	Per n-2 volte, calo di 1.
	Tot: 3/2 n-2 e sotto questo numero di confronti non si può andare.
	
	Osservazione:
	Dato una rray di lunghezza n, confronto ogni elemento a coppie di 2 a 2, per poi spostare in due gruppi rispettivamente i min e i max.
	Poi si cerca il minimo in min ed il massimo in max.
	Il problema è quando c'è un array dispari, una coppia sarà monca.
	
	-- Algoritmi per la ricerca del minimo
	Possibilmente con una complessità minore di nlogn. Possiamo usarne uno casuale.
	Supponiamo un array diviso in due parti [p-q], [q-r] t.c. gli alementi a sinistra siano <=, mentre a destra siano >=. Vogliamo l'indice i.
	
	Nel caso fortunato di un array diviso a metà avremo: T(n) = n + T(n/2). Per il teorema dell'esperto abbiamo che la parte nota vince sull'argomento, dunque la complessità è \Theta(n).
	Questo è tuttavia un caso fortunato. L'array non necessariamente è diviso a metà.
	Se tuttavia prendo a caso un elemento perno, verosimilmente, in media, avrò complessità lineare. E' una logica simile a randomized quicksort.
	
	RandSelect(A, p, r, i) {
		if (p = r) return A[p];
		
		q <- RandPartition(A, p, r);				// Divido l'array in due parti con un perno casuale.
		k <- q-p+1									// Quanti elementi stanno a sinistra del perno? (Quindi anche quanti stanno a dx)
		
		if (i <= k) return RandSelect(A,p,q,i);		// Seleziona l'iesimo elemento nell'array
		else return RandSelect(A, q+1, r, i-k);		// Seleziona dalla parte [q+1, r] e cerca l'indice i-k perché così non sballa gli indici del sub-array.
	}
	
	Complessità: T(n) = n + 1/n max(T(1), T(n-1)) + \sum_{i=1}^{n-1} max(T(i), T(n-i))
	
	Esiste anche un algoritmo deterministico per risolvre il problema delle selezioni in tempo lineare.
	Problema è che dipende molto dalle costanti, a tal punto che a grandi numero va meglio nlogn. Quindi non si può ragionare in termini di analisi asintotica.
	
	L'idea è pensare di scegliere il perno in maniera più intelligente, con lo scopo di stare in una zona proporzionale.
	Select del mediano non sarebbe male come idea, preferibilmente su un insieme più piccolo
	
	1: Dividi A in n/5 (approx per difetto) gruppi di 5 elementi più il gruppo con gli ultimi che restano.
	Compl: 0
	2: Calcolo il mediano di ogni singolo gruppo
	Compl: 5n
	3: Sia x il mediano dei mediani, calcolato ricorsivamente con select();
	Compl: T(n/5), con n/5 approx per eccesso. Eccesso a causa del gruppo con gli ultimi elementi restanti.
	4: Partiziona su x. Sia k il numero di elementi a sinistra. Quindi x è il perno usato.
	Compl: n
	5: Scelgo la parte interessata.
	if (i<=k) select(left, i)
	else select(right, i-k)
	Compl: T(dimensione più grande) [*qui]
	
	Complessità: \Theta(n) + T(n/5) + T(7/10n + 6)
	
	5 elementi perché così il mediano è ben definito, ed è anche un numero sufficientemente basso per avere complessità lineare. Un compromesso.
	
	Per esempio, con un array ottenuto dall'elaborazione dei mediani, ottengo che ha lunghezza 9. Se prendo la parte più grande, avrò 4 mediani maggiori del perno.
	Uno di questi mediani è in un insieme con possibilmente non 5 elementi. La cosa si mostra con:
	
	3 (\frac{n/5}{2} - 1 - 1) \\geq 3/10 n - 6		// Limite inferiore al numero di elementi che stanno a dx. Puoi farlo anche per la sx.
	
	Tutti gli elementi a destra meno uno esprimeranno 3 valori maggiori del mediano.
	
	n - (3/10 n -6) = 7/10n + 6. Questa formula è l'argomento da sostituire a "dimensione più grande"[*]
	
	L'equazione di ricorrenza \Theta(n) + T(n/5) + T(7/10n + 6) è sotto c(n). Usiamo la sostituzione
	
	T(n) \leq c'n + c(n/5) + c(7/10n + 6)		// Sostituzione applicabile solo se 7/10n + 6 è più piccolo di n.
	\leq c'n + c(n/5 + 1 + 7/10n + 6)
	= c'n + 7c + c9/10n					// Al prox passaggio procedo a tenere la mia ipotesi MENO all'espressione rimanente
	\leq cn - (1/10cn - c'n -7c)		// Serve (1/10cn - c'n -7c) \geq 0 per avere questo valore \leq cn.
\end{comment}